\documentclass[11pt]{article}
\usepackage[margin=1in]{geometry}
\usepackage{amsmath}
\usepackage{amssymb}
\usepackage{xcolor}

\title{Spring 2025 MC Problem 2: WARP Violation}
\author{ECO 310}
\date{}

\begin{document}

\maketitle

\section{The Problem}

On Monday, Homer Simpson buys a bundle of food and drink that costs \$90. On Tuesday, he buys a different bundle, for \$100, that would have cost \$80 at the Monday prices. In which of the following situations does he violate WARP:

\begin{enumerate}
    \item[(a)] If the Monday bundle would cost \$100 at the Tuesday prices.
    \item[(b)] If the Monday bundle would cost more than \$100 at the Tuesday prices.
    \item[(c)] Homer does not violate WARP.
\end{enumerate}

\textbf{Answer: (a)}

\section{Setting Up the Problem}

Let's define our notation clearly:
\begin{itemize}
    \item $p^M$ = Monday prices (price vector)
    \item $p^T$ = Tuesday prices (price vector)
    \item $x^M$ = Monday bundle (what Homer bought Monday)
    \item $x^T$ = Tuesday bundle (what Homer bought Tuesday)
\end{itemize}

\subsection{What We Know}

From the problem statement:
\begin{align}
p^M \cdot x^M &= 90 \tag{Monday bundle cost on Monday} \\
p^M \cdot x^T &= 80 \tag{Tuesday bundle cost at Monday prices} \\
p^T \cdot x^T &= 100 \tag{Tuesday bundle cost on Tuesday}
\end{align}

We need to determine what happens when $p^T \cdot x^M$ (Monday bundle cost at Tuesday prices) takes different values.

\section{Understanding WARP}

\textbf{Weak Axiom of Revealed Preference (WARP):} If bundle $x$ is chosen when bundle $y$ is affordable, then whenever $y$ is chosen, $x$ must not be affordable.

More formally: If $x$ is directly revealed preferred to $y$, then $y$ cannot be directly revealed preferred to $x$.

\subsection{Mathematical Statement}

If at prices $p$ and income $m$:
\[
p \cdot x \leq m \quad \text{and} \quad p \cdot y \leq m \quad \text{and we choose } x
\]

Then $x$ is \textbf{revealed preferred} to $y$ (denoted $x \succsim^R y$).

WARP says: If $x \succsim^R y$, then we cannot also have $y \succsim^R x$ (unless $x = y$).

\section{Applying WARP to Homer's Problem}

\subsection{Monday's Revealed Preference}

On Monday, Homer had budget \$90. He:
\begin{itemize}
    \item \textbf{Chose} bundle $x^M$ at cost $p^M \cdot x^M = 90$
    \item \textbf{Could afford} bundle $x^T$ at cost $p^M \cdot x^T = 80 < 90$
\end{itemize}

Since Homer chose $x^M$ when $x^T$ was strictly cheaper (and thus affordable), this reveals:
\[
\boxed{x^M \succ x^T} \quad \text{(Monday bundle is \textbf{strictly} preferred to Tuesday bundle)}
\]

This is a \textbf{strict} preference because he chose $x^M$ and didn't spend his full budget on $x^T$.

\subsection{Tuesday's Revealed Preference}

On Tuesday, Homer had budget \$100. He:
\begin{itemize}
    \item \textbf{Chose} bundle $x^T$ at cost $p^T \cdot x^T = 100$
    \item Bundle $x^M$ costs $p^T \cdot x^M = ?$
\end{itemize}

\textbf{Case (a): If $p^T \cdot x^M = 100$}

Homer could afford $x^M$ (costs exactly his budget), but chose $x^T$. This reveals:
\[
x^T \succsim x^M \quad \text{(Tuesday bundle is weakly preferred to Monday bundle)}
\]

\textbf{This is a WARP violation!} We have:
\begin{itemize}
    \item Monday revealed: $x^M \succ x^T$ (strict preference for Monday bundle)
    \item Tuesday revealed: $x^T \succsim x^M$ (weak preference for Tuesday bundle)
\end{itemize}

These preferences are \textbf{inconsistent}.

\textbf{Case (b): If $p^T \cdot x^M > 100$}

Homer \textbf{cannot afford} $x^M$ on Tuesday. Since $x^M$ is not in his budget set, we cannot infer anything about his preferences from Tuesday's choice. Therefore, \textbf{no WARP violation}.

\section{The WARP Violation Inequality}

Homer violates WARP when \textbf{both} of the following hold:

\begin{align}
p^M \cdot x^T &< p^M \cdot x^M \tag{Monday reveals $x^M \succ x^T$} \\
p^T \cdot x^M &\leq p^T \cdot x^T \tag{Tuesday reveals $x^T \succsim x^M$}
\end{align}

In our specific problem:
\begin{align}
80 &< 90 \quad \checkmark \tag{Monday condition satisfied} \\
p^T \cdot x^M &\leq 100 \tag{Tuesday condition - when does this hold?}
\end{align}

\textbf{Answer: WARP is violated when $p^T \cdot x^M \leq 100$}

This includes option (a) where $p^T \cdot x^M = 100$, but \textbf{not} option (b) where $p^T \cdot x^M > 100$.

\section{Intuition}

Think of it this way:
\begin{itemize}
    \item On Monday, Homer had a choice between the two bundles (both affordable) and picked Monday's bundle. This shows he prefers it.

    \item On Tuesday, if the Monday bundle is still affordable (costs $\leq$ \$100), but Homer chooses the Tuesday bundle instead, he's \textbf{contradicting} his Monday preference.

    \item However, if the Monday bundle becomes too expensive on Tuesday (costs $>$ \$100), then Homer simply can't afford it anymore. Choosing the Tuesday bundle doesn't contradict anything because he had no choice.
\end{itemize}

\section{General WARP Formula}

For any two observations $(p^1, x^1)$ and $(p^2, x^2)$, WARP is violated if:

\[
\boxed{
\begin{cases}
p^1 \cdot x^2 \leq p^1 \cdot x^1 & \text{(bundle 2 was affordable when we chose bundle 1)} \\
p^2 \cdot x^1 \leq p^2 \cdot x^2 & \text{(bundle 1 was affordable when we chose bundle 2)} \\
x^1 \neq x^2 & \text{(the bundles are different)}
\end{cases}
}
\]

with at least one inequality strict.

\section{Summary}

Homer violates WARP in option \textbf{(a)} because:
\begin{itemize}
    \item Monday: Chose $x^M$ over affordable $x^T$ $\Rightarrow$ $x^M \succ x^T$
    \item Tuesday: Chose $x^T$ over affordable $x^M$ $\Rightarrow$ $x^T \succsim x^M$
    \item These preferences contradict each other
\end{itemize}

Homer does \textbf{not} violate WARP in option (b) because $x^M$ is not affordable on Tuesday, so no preference reversal is revealed.

\end{document}
